
\chapter{Backend}
\label{chap:backend}

% Apresentar o backend do DoseCerta, sua responsabilidade
% na aplicação e sua relação com os demais componentes.

\section{Tecnologias e Ambiente de Desenvolvimento}
\label{sec:tecnologias-backend}

% Informar linguagem, framework, bibliotecas e ambiente utilizados.
% Justificar as escolhas em função dos requisitos do sistema.

\section{Arquitetura do Backend}
\label{sec:arquitetura-backend}

% Descrever a organização interna do backend:
% módulos, camadas e responsabilidades.
% Inserir diagrama de arquitetura, se disponível.

\section{API e Endpoints}
\label{sec:api-endpoints}

% Explicar como os clientes ou módulos do sistema acessam o backend.
% Apresentar os endpoints implementados, seus métodos HTTP,
% parâmetros, respostas e códigos de erro relevantes.

% Exemplo de tabela para documentar endpoints:
% \begin{table}[htbp]
%     \centering
%     \caption{Endpoints implementados no backend}
%     \label{tab:endpoints-backend}
%     \begin{tabular}{|l|l|l|p{6cm}|}
%         \hline
%         Método & Rota & Função & Descrição \\
%         \hline
%         % GET & /exemplo & Consultar & Descrever a operação \\
%         \hline
%     \end{tabular}
% \end{table}

\section{Modelo de Dados e Persistência}
\label{sec:modelo-dados-backend}

% Descrever os dados armazenados e sua organização.
% Apresentar entidades, atributos e relacionamentos.
% Incluir diagrama entidade-relacionamento ou esquema,
% caso exista banco de dados.

\section{Regras de Negócio}
\label{sec:regras-negocio-backend}

% Descrever as regras implementadas no servidor,
% como cadastro, configuração de medicamentos ou horários,
% conforme as funcionalidades existentes.
% Diferenciar validações do backend das validações da interface.

\section{Autenticação e Controle de Acesso}
\label{sec:autenticacao-backend}

% Descrever o mecanismo de autenticação e autorização,
% se implementado. Informar como credenciais ou tokens são
% tratados, sem expor senhas, chaves ou segredos.

\section{Validação e Tratamento de Erros}
\label{sec:validacao-erros-backend}

% Explicar a validação dos dados recebidos,
% o tratamento de requisições inválidas e os erros retornados.
% Registrar medidas de segurança efetivamente implementadas.

\section{Integração com os Demais Componentes}
\label{sec:integracao-backend}

% Descrever a comunicação do backend com o aplicativo,
% sistema embarcado ou outros serviços.
% Informar formatos de dados e protocolos utilizados,
% conforme a implementação.

\section{Implantação e Execução}
\label{sec:implantacao-backend}

% Informar como o backend é executado no ambiente de desenvolvimento
% ou implantação: configuração, dependências e serviços necessários.
% Não incluir credenciais ou informações sensíveis.

\section{Síntese do Backend}
\label{sec:sintese-backend}

% Resumir as funcionalidades implementadas,
% o papel do backend na solução e suas limitações.
